\section{Discussion and limitations}
\label{sec:discussion}
\paragraph{Significance for representation learning.}
The central object is the attainable risk of a declared representation--decoder pair. A risk lower certificate supports a conclusion that a fitted-model comparison cannot: every member fails the specified requirement. A constructive class upper witness can support the opposite diagnosis even when evaluated fits fail. The centered hierarchy supplies increasingly complete population lower information; the empirical studies show when that information changes a justified conclusion and when it cannot.

\paragraph{What the hierarchy does not identify.}
The theorem concerns bounded scalar targets, a fixed representation metric, and a fixed sensitivity budget. Compact support and unconditional iid sampling enter the consistency statement. A common-origin conditional construction can converge to an average of conditional optima instead of the optimum of one shared decoder. Vector-output and arbitrary dependent-series consistency do not follow. A positive lower endpoint can arise from irreducible uncertainty or decoder restrictions as well as representation information loss.

\paragraph{Three distinct sources of inconclusiveness.}
A finite observation block can relax too many cross-observation constraints; a valid statistical endpoint can lack power; and a numerical lower witness can be loose. Increasing the optimization effort cannot remove an impossible calibration floor. Increasing block size can reduce population bias while worsening statistical resolution. When a class already has a constructive adequate witness on a cohort, no valid class-exclusion lower bound can cross the same threshold.

\paragraph{Scope of representation and cost claims.}
Class nesting preserves feasible functions, not the quality of a particular regularized fit. Exact reparameterization can preserve information while changing a Euclidean Lipschitz class or the ease of learning within it. The selected KAHM preprocessing and caps remain task dependent. Strong raw, kernel and conventional predictors prevent the certificate construction from being mistaken for a general accuracy advantage. Minimum feasible cost is relative to a declared family and cost functional; feature count alone does not certify end-to-end latency or storage.

\paragraph{Sampling and adaptation.}
Independent calibration units and the declared confidence ledgers support the synthetic studies' risk statements; repeated learning runs describe procedure variability. A deterministic census and chronological benchmark periods do not inherit those population guarantees. Nearby windows, fitting seeds on the same targets, and algebraically related responses are not independent population replications. Development selection must precede independent calibration, or its additional adaptivity must be covered. Reusing exposed cohorts is valuable for diagnosis but remains labeled as such. No campaign-wide significance claim is obtained by counting completed studies.

\paragraph{Numerical assurance.}
Weak duality and constructive feasibility can be checked independently of a numerical solver's termination status. Exact arithmetic still depends on the chosen function semantics and trusted inputs. Native feature enclosures demonstrated on bounded manufactured configurations do not automatically certify every feature computation in the public-data studies. The later affine output checker assumes trusted stored coordinates, recomputes an exact reference, and has measurable overhead. It provides neither cryptographic authentication nor a speedup over direct exact recomputation.

\paragraph{Further applications.}
A public-data population application needs a defensible sampling or dependence model and an admissible predictor under the same loss and class. Certificate-directed savings from avoiding fits, reliable adaptive discovery of refinements, and end-to-end native feature computation certification remain separate extensions. They are not premises of the scoped results presented here.

\section{Conclusion}
Prediction-risk certification should distinguish the performance of a fitted decoder from the best risk allowed by its representation and sensitivity class. The centered scalar hierarchy connects a pairwise obstruction to that optimum under explicit conditions. Valid calibration, verified witnesses, and fitted-risk upper bounds support interpretable exclusion, adequacy, and family-relative cost decisions. The studies demonstrate both informative certificates and limits on their practical use. Separating population risk, finite-cohort capacity, and numerical output fidelity makes those conclusions usable without overstating generality.

